(1) 简支两端的空间正弦模满足位移及弯矩为零的边界：
\[
X_n(x)=\sin(n\pi x/L),\quad k_n=n\pi/L,\quad n=1,2,\dots.
\]
因为 \(X_n''=-k_n^2X_n\)、\(X_n''''=k_n^4X_n\)，令 \(y=X_n e^{i\omega_nt}\) 代入得
\[
\omega_n^2(P)=\frac{EI\,k_n^4-Pk_n^2}{\rho A}
=\frac{k_n^2(EI\,k_n^2-P)}{\rho A}.
\]
在质量权重 \(\rho A\) 下，不同 \(n\) 的正弦函数彼此正交。

(2) 第 \(n\) 模稳定需要 \(\omega_n^2>0\)，即 \(P<EI k_n^2\)。全部模均稳定的最严格条件来自最低频的 \(n=1\)：
\[
0\le P<P_{\rm cr},\qquad P_{\rm cr}=\frac{\pi^2EI}{L^2}.
\]
在 \(P=P_{\rm cr}\) 时第一模 \(\omega_1=0\)，为中性静态弯曲；当压力超过该值，第一模的 \(\omega_1^2<0\)，出现指数增长，直梁小扰动在此线性化模型中不稳定。

(3) 对 PDE 乘以 \(y_t\) 并沿全梁积分，分部积分，简支边界使边界项消失，得到
\[
\frac{d}{dt}\mathcal E=0,\qquad
\mathcal E=\frac12\int_0^L\left[\rho A y_t^2+EI y_{xx}^2-P y_x^2\right]dx.
\]
弹性弯曲能提供 \(+EI y_{xx}^2/2\)，外部压缩死载荷使位移的二阶有效势能包含负的 \(-P y_x^2/2\)。代入正弦级数，每个系数对应势能因子 \(EI k_n^4-Pk_n^2\)；恰在 \(P_{\rm cr}\) 时最低模势能曲率为零。负势能二次型并不违反能量守恒，但可导致动态失稳。

(4) 对最低模
\[
\omega_1=\frac{\pi}{L\sqrt{\rho A}}\sqrt{P_{\rm cr}-P}
\quad (0\le P<P_{\rm cr}),
\]
所以临界附近 \(\omega_1\propto(P_{\rm cr}-P)^{1/2}\)，振动周期发散。若 \(P>P_{\rm cr}\)，令 \(\sigma=\pi\sqrt{P-P_{\rm cr}}/(L\sqrt{\rho A})\)，第一模时间因子为 \(e^{\pm\sigma t}\)，其中增长分支揭示直梁失稳。但随着挠度增大，几何非线性、边界大转角和材料效应必须纳入；仅靠线性 PDE 不能计算有限幅值的后屈曲平衡。
